\exercisesection{\S~4}

本段中所考察的李代数不一定是有限维的。

\begin{enumerate}
\def\labelenumi{\arabic{enumi})}
\tightlist
\item
  保留第2节的记号。令 \(\lambda \in k^{\mathrm{B}}\)。对任意 \(\alpha \in \mathrm{B}\)，将向量空间 E 的自同态 \(X_{-\alpha}^{\lambda}, H_{\alpha}^{\lambda}, X_{\alpha}^{\lambda}\) 与之关联，使得
\end{enumerate}

\[
X _ {- \alpha} ^ {\lambda} (\alpha_ {1}, \ldots , \alpha_ {n}) = (\alpha , \alpha_ {1}, \ldots , \alpha_ {n})
\]

\[
H _ {\alpha} ^ {\lambda} (\alpha_ {1}, \ldots , \alpha_ {n}) = (\lambda (\alpha) - \sum_ {i = 1} ^ {n} n (\alpha_ {i}, \alpha)) (\alpha_ {1}, \ldots , \alpha_ {n}).
\]

向量 \(X_{\alpha}^{\lambda}(\alpha_{1},\ldots,\alpha_{n})\) 由关于 n 的归纳用下式定义

\[
X _ {\alpha} ^ {\lambda} (\alpha_ {1}, \ldots , \alpha_ {n}) = (X _ {- \alpha_ {1}} ^ {\alpha} X _ {\alpha} ^ {\lambda} - \delta_ {\alpha , \alpha_ {1}} H _ {\alpha} ^ {\lambda}) (\alpha_ {2}, \ldots , \alpha_ {n}),
\]

其中 \(\delta_{\alpha, \alpha_1}\) 是 Kronecker 符号；约定，如果 \((\alpha_1, \ldots, \alpha_n)\) 是空字，则 \(X_{\alpha}^{\lambda}(\alpha_1, \ldots, \alpha_n)\) 为零。

证明引理1和2对这些自同态仍然成立。于是得到一个表示 \(\rho_{\lambda}:\mathfrak{a}\to \mathfrak{gl}(\mathrm{E})\)，使得

\[
\rho_ {\lambda} (x _ {\alpha}) = X _ {\alpha} ^ {\lambda}, \rho_ {\lambda} (h _ {\alpha}) = H _ {\alpha} ^ {\lambda}, \rho_ {\lambda} (x _ {- \alpha}) = X _ {- \alpha} ^ {\lambda}.
\]

¶ 2) 保留第3节的记号。设 m 是 a 的理想。

\begin{enumerate}
\def\labelenumi{\alph{enumi})}
\item
  令 \(\alpha\) 和 \(\beta\) 为 B 中两个不同的元素。假设当 N 充分大时，(ad \(x_{\alpha})^{\mathrm{N}}x_{\beta}\) 属于 \(\mathfrak{m}\)。证明 \(x_{\alpha \beta} \in \mathfrak{m}\)。（将§1第2节的结果应用于 \(\mathfrak{a} / \mathfrak{m}\)，其中赋予适当的 \(\mathfrak{sl}(2,k)\)-模结构。）
\item
  证明 \(n + \theta n\) 是 a 的最小有限余维理想。证明它也是包含 \((\operatorname{ad} x_{\alpha})^{4}x_{\beta}\) 和 \((\operatorname{ad} x_{-\alpha})^{4}x_{-\beta}\) 的最小理想。
\end{enumerate}

\begin{enumerate}
\def\labelenumi{\arabic{enumi})}
\setcounter{enumi}{2}
\tightlist
\item
  令 \((\mathfrak{g},\mathfrak{h})\) 为分裂半单李代数，R 为相应的根系，B 为 R 的一个基。对任意 \(\alpha \in \mathrm{B}\)（分别地，对任意一对 \((\alpha ,\beta)\in \mathrm{B}^2\)），令 \(\mathrm{R}(\alpha)\)（分别地，令 \(\mathrm{R}(\alpha ,\beta)\)）为 R 中由 \(\pm \alpha\) 构成的闭子集（分别地，为包含 \(\pm \alpha\) 和 \(\pm \beta\) 的最小闭子集）。令 \(\mathfrak{g}(\alpha)\)（分别地，令 \(\mathfrak{g}(\alpha ,\beta)\)）为代数 \(\mathfrak{h} + \mathfrak{g}^{\mathrm{R}(\alpha)}\)（分别地，代数 \(\mathfrak{h} + \mathfrak{g}^{\mathrm{R}(\alpha ,\beta)}\)）的导代数，参见§3。
\end{enumerate}

\begin{enumerate}
\def\labelenumi{\alph{enumi})}
\item
  证明 \(\mathfrak{g}(\alpha) = kH_{\alpha}\oplus \mathfrak{g}^{\alpha}\oplus \mathfrak{g}^{-\alpha}\)；它同构于 \(\mathfrak{sl}(2,k)\)。
\item
  证明 \(\mathfrak{g}(\alpha,\beta)\) 是半单的，且由 \(\mathfrak{g}(\alpha)\) 和 \(\mathfrak{g}(\beta)\) 生成。其根系可以与 \(\mathrm{R}(\alpha,\beta)\) 相识别。
\item
  令 \(\mathfrak{s}\) 为李代数（不一定有限维）。对所有 \(\alpha \in \mathrm{B}\)，令 \(f_{\alpha}\) 为从 \(\mathfrak{g}(\alpha)\) 到 \(\mathfrak{s}\) 的同态。假设对任意一对 \((\alpha, \beta)\)，存在同态 \(f_{\alpha \beta}: \mathfrak{g}(\alpha, \beta) \to \mathfrak{s}\)，同时延拓 \(f_{\alpha}\) 和 \(f_{\beta}\)。证明在这种情况下，存在唯一同态 \(f: \mathfrak{g} \to \mathfrak{s}\)，延拓各个 \(f_{\alpha}\)。（使用命题4 (i)。）
\end{enumerate}

\begin{enumerate}
\def\labelenumi{\arabic{enumi})}
\setcounter{enumi}{3}
\item
  令 \(\mathfrak{g}\) 为可分裂半单李代数，令 \(\sigma\) 为 \(k\) 的自同构。令 \(\mathfrak{g}_{\sigma}\) 为由 \(\mathfrak{g}\) 借助 \(\sigma\) 扩张标量所得到的李代数。证明 \(\mathfrak{g}_{\sigma}\) 同构于 \(\mathfrak{g}\)。（使用命题4的推论。）
\item
  \begin{enumerate}
  \def\labelenumii{\alph{enumii})}
  \tightlist
  \item
    令 g 为单李代数，令 \(k_{1}\) 为 g 的伴随表示的交换子。证明 \(k_{1}\) 是交换域，是 k 的有限扩张，并且 g 是绝对单的 \(k_{1}\)-李代数。
  \end{enumerate}
\end{enumerate}

反之，若 \(k_{1}\) 是 k 的有限扩张，且 g 是绝对单的 \(k_{1}\)-李代数，则 g 是单的 k-李代数，并且其伴随表示的交换子可以与 \(k_{1}\) 相识别。

\begin{enumerate}
\def\labelenumi{\alph{enumi})}
\setcounter{enumi}{1}
\tightlist
\item
  令 \(k'\) 为包含 \(k_{1}\) 的 Galois 扩张。证明 \(\mathfrak{g}_{(k')}\) 是 \([k_{1}:k]\) 个绝对单代数的直积。当 \(\mathfrak{g}_{(k')}\) 分裂时，这些代数彼此同构。（使用习题4。）
\end{enumerate}

\begin{enumerate}
\def\labelenumi{\arabic{enumi})}
\setcounter{enumi}{5}
\tightlist
\item
  令 \(\mathbf{A}\) 为交换环，令 \(\mathfrak{u}\) 为由生成族 \(\{x, y\}\) 及下列关系定义的 A-李代数
\end{enumerate}

\[
[ x, [ x, y ] ] = 0, \quad [ y, [ y, [ y, x ] ] ] = 0.
\]

证明 \(\mathfrak{u}\) 是以

\[
\{x, y, [ x, y ], [ y, [ x, y ] ] \}.
\]

证明当 \(\mathrm{A} = k\) 时，\(\mathfrak{u}\) 同构于 \(\mathfrak{a}_{+} / \mathfrak{n}\)，它与 \(\mathrm{B}_2\) 型根系相应。

¶7) 令 A 为其中 2 可逆的交换环，令 u 为由生成族 \(\{x,y\}\) 及下列关系定义的 A-李代数

\[
[ x, [ x, y ] ] = 0, \quad [ y, [ y, [ y, [ y, x ] ] ] ] = 0.
\]

证明 \(\mathfrak{u}\) 是以

\[
\left. \right.\left\{x, y, [ x, y ], [ y, [ x, y ] ], [ y, [ y, [ x, y ] ] ], [ x, [ y, [ y, [ x, y ] ] ] ] \right\}.
\]

证明当 A = k 时，u 同构于代数 \(a_{+}/n\)，它与 \(G_{2}\) 型根系相应。

